\chapter{Cohomology}
\label{pa-ch-cohomology}

\begin{theorem}[Finiteness and base change]
  \label{pa-thm-cohomology}
  \dcref{III.12,custom}
  \uses{pa-def-genus}
  \cite{stacks-project}
  If \(\mathcal F\) is coherent on \(C\), then \(H^i(C,\mathcal F)\) is finite
  dimensional over \(k\), and for every extension \(k\subset k'\),
  \[
    H^i(C,\mathcal F)\otimes_k k'\cong
    H^i(C_{k'},\mathcal F_{k'}).
  \]
\end{theorem}

\begin{proof}
  A finite affine cover of the proper curve gives a bounded Cech complex of
  finite-dimensional \(k\)-vector spaces computing coherent cohomology.
  Extending scalars tensors this complex with \(k'\); flatness of a field
  extension commutes with its cohomology.  This proves both finiteness and the
  displayed base-change isomorphism.
\end{proof}

\begin{proposition}[Global functions]
  \label{pa-prop-global-functions}
  \dcref{III.1,custom}
  \cite{hartshorne-ag,stacks-project}
  Geometric integrality and properness imply \(H^0(C,\mathcal O_C)=k\).
\end{proposition}

\begin{proof}
  Apply base change to \(H^0\) and an algebraic closure \(\bar k\).  The
  proper integral scheme \(C_{\bar k}\) has only constant global functions,
  so
  \(H^0(C,\mathcal O_C)\otimes_k\bar k\simeq\bar k\).  The structural
  inclusion of \(k\) is consequently an isomorphism.
\end{proof}

\begin{proposition}[Picard tangent space]
  \label{pa-prop-picard-tangent}
  \dcref{5.11}
  \uses{pa-def-genus,pa-thm-cohomology}
  \cite[Thm.~5.11]{kleiman-picard}
  For the identity element of the relative Picard scheme,
  \[
    T_0\operatorname{Pic}_{C/k}\cong H^1(C,\mathcal O_C),
    \qquad \dim_kT_0\operatorname{Pic}_{C/k}=g(C).
  \]
\end{proposition}

\begin{proof}
  Put \(k[\epsilon]=k[\epsilon]/(\epsilon^2)\).  The kernel of the reduction
  map on the etale Picard functor is computed by
  \[
    1\longrightarrow1+\epsilon\mathcal O_C
    \longrightarrow\mathcal O_{C_{k[\epsilon]}}^\times
    \longrightarrow\mathcal O_C^\times\longrightarrow1.
  \]
  The first term is additively \(\mathcal O_C\), so the kernel is
  \(H^1(C,\mathcal O_C)\).  Scalar substitutions
  \(\epsilon\mapsto a\epsilon\) supply the vector-space structure and the
  preceding base-change theorem gives its compatibility with \(k'\).
\end{proof}
